% Generated from the reviewed LONGFORM_GENERATED_V1 complete reader input. Do not hand-edit.
\documentclass[lualatex,a4paper,10pt]{jlreq}
\usepackage{jgaisurvey}
\addbibresource{references.bib}

\surveysetup{SP001}{Japanese Generative AI Technical Survey Special}{Thematic history / as of 2026-08-21 17:00:00 UTC}{As-of boundary: 2026-08-21 17:00:00 UTC}
\surveyeditiondescriptor{Thematic Longform Special}
\surveycoverstory{Synthetic longform headline (thematic-full-cuu-rel-v2)}{Synthetic longform deck (thematic-full-cuu-rel-v2): revised lineage carried.}{synthetic-longform-anchor-thematic-full-cuu-rel-v2}
\surveyrepository{eariver/japanese-generative-ai-survey}
\surveybuildvalue{LuaLaTeX / jlreq / LuaTeX-ja}
\surveycoverlabels{Editorial cutoff}{Repository}{Build}
\surveystrapline{一次情報・論文・OSS・X上の技術反応を分離し、検証可能なEvidence chainとして編成する}
\renewcommand{\contentsname}{目次}

\begin{document}
\surveycover
\clearpage
\section*{Synthetic longform frontmatter (thematic-full-cuu-rel-v2)}
\addcontentsline{toc}{section}{Synthetic longform frontmatter (thematic-full-cuu-rel-v2)}
Synthetic lede (thematic-full-cuu-rel-v2): independent reasoning on the revised graph.

\begin{claimboundary}[Scope / attribution boundary]
\noindent Synthetic longform scope note (thematic-full-cuu-rel-v2)\par
\end{claimboundary}
\medskip
\tableofcontents
\clearpage
% package:pkg-core
\Needspace{0.24\textheight}
\section{Synthetic headline for pkg-core (thematic-full-cuu-rel-v2): revised MATERIAL finding carried}
\label{pkg:pkg-core}
\sectionkicker{THEMATIC LINEAGE 1/1 — core}
\noindent\textbf{Synthetic deck for pkg-core (thematic-full-cuu-rel-v2): independent reasoning on the revised graph.} \cite{sp001lnd001}\par\medskip
\begin{themeoverview}[Theme at a glance]
\begin{tabularx}{\linewidth}{@{}>{\bfseries}p{0.22\linewidth}X@{}}
Glance alpha (thematic-full-cuu-rel-v2) & Synthetic glance alpha prose (thematic-full-cuu-rel-v2). \cite{sp001lnd001} \\
Glance beta (thematic-full-cuu-rel-v2) & Synthetic glance beta prose (thematic-full-cuu-rel-v2). \cite{sp001lnd001} \\
\end{tabularx}
\end{themeoverview}
\medskip
\begin{multicols}{2}
\subsection*{Synthetic narrative one (thematic-full-cuu-rel-v2)}
\noindent Synthetic narrative one prose (thematic-full-cuu-rel-v2). \cite{sp001lnd001}\par\medskip
\noindent Synthetic narrative one more prose (thematic-full-cuu-rel-v2). \cite{sp001lnd001}\par\medskip
\subsection*{Synthetic narrative two (thematic-full-cuu-rel-v2)}
\noindent Synthetic narrative two prose (thematic-full-cuu-rel-v2). \cite{sp001lnd001}\par\medskip
\noindent Synthetic narrative two more prose (thematic-full-cuu-rel-v2). \cite{sp001lnd001}\par\medskip
\end{multicols}
\medskip
\begin{wideflow}
\subsection*{Transition timeline}
\begin{tabularx}{\linewidth}{@{}>{\bfseries}p{0.20\linewidth}X@{}}
Timeline alpha (thematic-full-cuu-rel-v2) & Synthetic timeline alpha prose (thematic-full-cuu-rel-v2). \cite{sp001lnd001} \\
Timeline beta (thematic-full-cuu-rel-v2) & Synthetic timeline beta prose (thematic-full-cuu-rel-v2). \cite{sp001lnd001} \\
\end{tabularx}
\end{wideflow}
\medskip
\begin{multicols}{2}
\subsection*{Synthetic package synthesis (thematic-full-cuu-rel-v2)}
\noindent Synthetic synthesis prose (thematic-full-cuu-rel-v2). \cite{sp001lnd001}\par\medskip
\end{multicols}
\medskip
\begin{claimboundary}[Claim boundary]
\noindent Synthetic boundary prose (thematic-full-cuu-rel-v2). \cite{sp001lnd001}\par
\end{claimboundary}
\medskip
\subsection*{Source-backed Technical Notes}
\begin{technicalnote}[Synthetic note (thematic-full-cuu-rel-v2)]
\noindent\textbf{Chronology:} Synthetic chronology (thematic-full-cuu-rel-v2). \cite{sp001lnd001}\par
\smallskip\noindent\textbf{Technical points:}\par
\begin{itemize}
\item Synthetic point one (thematic-full-cuu-rel-v2). \cite{sp001lnd001}
\item Synthetic point two (thematic-full-cuu-rel-v2). \cite{sp001lnd001}
\end{itemize}
\noindent\textbf{Limitation / attribution:} Synthetic limitation (thematic-full-cuu-rel-v2).\par
\smallskip\noindent\textbf{Primary URL:} \url{https://example.invalid/synthetic-lineage-model}\par
\end{technicalnote}
\smallskip
\Needspace{0.30\textheight}
\section{Synthetic cross-family heading (thematic-full-cuu-rel-v2)}
\sectionkicker{CROSS-FAMILY SYNTHESIS}
\noindent Synthetic cross prose alpha (thematic-full-cuu-rel-v2). \cite{sp001lnd001}\par\medskip
\begin{themeoverview}[Cross-family comparison]
\scriptsize
\begin{tabularx}{\linewidth}{@{}>{\bfseries}p{0.12\linewidth}XXXX@{}}
\toprule
観点 & GLM & Qwen & DeepSeek & Kimi \\ 
\midrule
Synthetic dimension a (thematic-full-cuu-rel-v2) & Synthetic glm note a (thematic-full-cuu-rel-v2). & Synthetic qwen note a (thematic-full-cuu-rel-v2). & Synthetic deepseek note a (thematic-full-cuu-rel-v2). & Synthetic kimi note a (thematic-full-cuu-rel-v2). \cite{sp001lnd001} \\
Synthetic dimension b (thematic-full-cuu-rel-v2) & Synthetic glm note b (thematic-full-cuu-rel-v2). & Synthetic qwen note b (thematic-full-cuu-rel-v2). & Synthetic deepseek note b (thematic-full-cuu-rel-v2). & Synthetic kimi note b (thematic-full-cuu-rel-v2). \cite{sp001lnd001} \\
Synthetic dimension c (thematic-full-cuu-rel-v2) & Synthetic glm note c (thematic-full-cuu-rel-v2). & Synthetic qwen note c (thematic-full-cuu-rel-v2). & Synthetic deepseek note c (thematic-full-cuu-rel-v2). & Synthetic kimi note c (thematic-full-cuu-rel-v2). \cite{sp001lnd001} \\
\bottomrule
\end{tabularx}
\end{themeoverview}
\Needspace{0.28\textheight}
\section{この号の総括}
\label{sec:issue-summary}
\sectionkicker{ISSUE SYNTHESIS}
\begin{multicols}{2}
\noindent Synthetic closing paragraph alpha (thematic-full-cuu-rel-v2).\par\medskip
\noindent Synthetic closing paragraph beta (thematic-full-cuu-rel-v2).\par\medskip
\noindent Synthetic closing paragraph gamma (thematic-full-cuu-rel-v2).\par\medskip
\end{multicols}
\clearpage
\onecolumn
\printbibliography[title={References / Source Notes}]
\end{document}
